\documentclass[multi=tikzpicture,border=8pt]{standalone}
\usepackage[UTF8]{ctex}
\usepackage{amsmath,amssymb,tikz}
\usetikzlibrary{arrows.meta,calc,decorations.pathmorphing}
\definecolor{blue}{HTML}{17365D}
\definecolor{red}{HTML}{AC3931}
\tikzset{every node/.style={font=\small},arrow/.style={-{Stealth[length=2.8mm]},thick},dim/.style={|<->|,thin},box/.style={draw,rounded corners,align=center,minimum width=2.5cm,minimum height=.9cm}}
\begin{document}
\begin{tikzpicture}

\draw[thick,blue](0,-.9)rectangle(5,.9);\draw[thick,red](3,-1.2)rectangle(8,1.2);
\foreach \x in {.2,.45,...,4.9}{\draw[blue](\x,-.9)--(\x,.9);}\foreach \x in {3.2,3.45,...,7.9}{\draw[red](\x,-1.2)--(\x,1.2);}
\draw[arrow](1,0)--(2.4,0)node[above,midway]{$B_0$};\draw[arrow,red](5,1.3)--(5,2.3)node[above]{$B_r$};\draw[dim](3,-1.7)--(5,-1.7)node[below,midway]{重叠$q$};\node[above]at(1.5,1.4){管1：$n,I$};\node[above]at(6.5,1.6){管2：$n,I$};

\end{tikzpicture}
\begin{tikzpicture}

\draw[thick](0,-1)--(5,-1);\draw[thick](0,1)--(5,1);\draw[thick](0,0)ellipse(.35 and 1);\draw[thick](5,0)ellipse(.35 and 1);
\draw[arrow](-.1,0)--(-1.5,0)node[left]{$-SB_x$};\draw[arrow](5.1,0)--(6.5,0)node[right]{$S(B_x+dB_x)$};\draw[arrow,red](2.5,1)--(2.5,2)node[above]{$2\pi rB_rdx$};
\draw[dim](0,-1.6)--(5,-1.6)node[midway,below]{$dx$};\node[below]at(2.5,-2.5){端面净通量＋侧面通量$=0$};

\end{tikzpicture}
\begin{tikzpicture}

\draw[blue,thick](0,0)circle(2);\foreach \a in {0,90,180,270}{\draw[arrow,red]({1.3*cos(\a)},{1.3*sin(\a)})--({2.5*cos(\a)},{2.5*sin(\a)});\draw[arrow,blue]({2*cos(\a-20)},{2*sin(\a-20)})arc(\a-20:\a+12:2);}
\node[red]at(3.3,0){$B_r>0$};\node[blue]at(-3.1,1){$I$逆时针};\node[align=center]at(0,0){$\otimes F_x$\\轴力向纸内};\node[below]at(0,-2.7){以$+x$朝向观察者，$\mathbf e_\varphi\times\mathbf e_r=-\mathbf e_x$};

\end{tikzpicture}
\begin{tikzpicture}
\draw[arrow](0,0)--(7.4,0)node[right]{$x/a$};\draw[arrow](0,0)--(0,3.4)node[above]{$B_x/B_0$};
\draw[blue,very thick,domain=0:7,samples=90]plot(\x,{1.5*(1+tanh(\x-3.5))});\draw[red,dashed,thick,domain=0:7,samples=90]plot(\x,{1.5*(1+tanh(2*(\x-3.5)))});
\node[left]at(0,3){1};\node[left]at(0,0){0};\node[align=center]at(3.5,-1.1){不同过渡宽度\\$\int dB_x=B_0-0$};

\end{tikzpicture}
\begin{tikzpicture}

\node[box,fill=red!8](a)at(0,0){两电源\\$2I_1I_2dM$};\node[box,fill=blue!8](b)at(5,1.4){磁场能增加\\$I_1I_2dM$};\node[box](c)at(5,-1.4){机械功\\$F_qdq=I_1I_2dM$};\draw[arrow](a)--(b);\draw[arrow](a)--(c);\node[below]at(2.5,-2.8){保持电流不变，电源持续交换能量};

\end{tikzpicture}
\end{document}
